\chapter{Express Routing}

\section{HTTP Methods and Their Purpose}

\begin{itemize}
  \item \textbf{GET} --- retrieve a resource or collection; no body, safe and idempotent.
  \item \textbf{POST} --- create a new resource; carries a request body.
  \item \textbf{PUT} --- replace a resource entirely (send the full object).
  \item \textbf{PATCH} --- partially update a resource (send only changed fields).
  \item \textbf{DELETE} --- remove a resource.
\end{itemize}

\begin{notebox}[NOTE]
In practice many REST APIs use \texttt{PUT} where \texttt{PATCH} would be more semantically correct. This handbook uses \texttt{PUT} for task updates since the frontend form always sends the full task object.
\end{notebox}

\section{Route Files with express.Router()}

\whatwhyhow{
\texttt{express.Router()} creates a mini, mountable Express application scoped to one resource.
}{
Without it, every route would live in one giant \texttt{app.js}, making the file unmanageable as the API grows.
}{
Define routes in \texttt{routes/*.routes.js}, export the router, and mount it in \texttt{app.js} with \texttt{app.use("/api/tasks", taskRoutes)}.
}

\subsection{routes/auth.routes.js}

\begin{lstlisting}[language=JavaScript]
const express = require("express");
const { register, login, logout, getMe } = require("../controllers/auth.controller");
const { protect } = require("../middleware/auth.middleware");

const router = express.Router();

router.post("/register", register);
router.post("/login", login);
router.post("/logout", logout);
router.get("/me", protect, getMe);

module.exports = router;
\end{lstlisting}

\subsection{routes/user.routes.js}

\begin{lstlisting}[language=JavaScript]
const express = require("express");
const { getUsers, getUserById, updateUser } = require("../controllers/user.controller");
const { protect } = require("../middleware/auth.middleware");

const router = express.Router();

router.get("/", protect, getUsers);
router.get("/:id", protect, getUserById);
router.put("/:id", protect, updateUser);

module.exports = router;
\end{lstlisting}

\subsection{routes/task.routes.js}

\begin{lstlisting}[language=JavaScript]
const express = require("express");
const {
  createTask,
  getTasks,
  getTaskById,
  updateTask,
  deleteTask,
} = require("../controllers/task.controller");
const { protect } = require("../middleware/auth.middleware");

const router = express.Router();

router.use(protect); // every task route requires authentication

router.post("/", createTask);
router.get("/", getTasks);
router.get("/:id", getTaskById);
router.put("/:id", updateTask);
router.delete("/:id", deleteTask);

module.exports = router;
\end{lstlisting}

\subsection{routes/profile.routes.js}

\begin{lstlisting}[language=JavaScript]
const express = require("express");
const { getProfile, updateProfile } = require("../controllers/profile.controller");
const { protect } = require("../middleware/auth.middleware");

const router = express.Router();

router.get("/", protect, getProfile);
router.put("/", protect, updateProfile);

module.exports = router;
\end{lstlisting}

\section{Route Params vs Query Params vs Body}

\begin{itemize}
  \item \textbf{Route params} --- part of the URL path, used to identify a specific resource. \texttt{/tasks/:id} $\rightarrow$ \texttt{req.params.id}.
  \item \textbf{Query params} --- key/value pairs after \texttt{?}, used for filtering, pagination, sorting. \texttt{/tasks?page=1\&limit=10} $\rightarrow$ \texttt{req.query.page}, \texttt{req.query.limit}.
  \item \textbf{Request body} --- JSON payload sent with POST/PUT/PATCH, used to create or update data. \texttt{req.body} is only populated because \texttt{express.json()} middleware ran first.
\end{itemize}

\begin{lstlisting}[language=JavaScript]
// GET /api/tasks?page=2&limit=10&status=todo
const getTasks = async (req, res) => {
  const { page = 1, limit = 10, status } = req.query;
  // ...build filter and pagination from query params
};

// GET /api/tasks/:id
const getTaskById = async (req, res) => {
  const { id } = req.params; // route param identifies the exact task
};
\end{lstlisting}

\section{Task Routes Reference}

\begin{table}[h]
\centering
\begin{tabularx}{\textwidth}{l l X}
\toprule
\textbf{Method} & \textbf{Path} & \textbf{Purpose} \\
\midrule
POST   & /api/tasks      & Create a new task for the logged-in user \\
GET    & /api/tasks      & List tasks (supports pagination \& filtering via query params) \\
GET    & /api/tasks/:id  & Fetch a single task by its ObjectId \\
PUT    & /api/tasks/:id  & Replace/update an existing task \\
DELETE & /api/tasks/:id  & Remove a task \\
\bottomrule
\end{tabularx}
\end{table}

\begin{tipbox}[TIP]
Keep route files thin. They should only wire paths to controller functions and attach middleware — no business logic, no direct database calls.
\end{tipbox}
